\chapter{Фазовый синхронизм при фотоиндуцированной генерации второй гармоники в оптическом волокне}\label{ch:nonlin2}

\section{Постановка задачи}


Несмотря на многолетнее изучение генерации второй гармоники (ВГГ) в германосиликатных волокнах, динамика формирования решетки квадратичной восприимчивости в процессе оптического полинга исследована неполно. В частности, эволюция спектральных и температурных кривых настройки ВГГ, являющихся наиболее чувствительным индикатором изменения параметров решетки и условий фазового синхронизма, до настоящего времени экспериментально не изучалась. Существующие теоретические модели требуют верификации, а понимание этих процессов необходимо как для создания эффективных источников на основе ВГГ, так и для подавления паразитной генерации в волоконных лазерных системах.

Целью данной работы является экспериментальное исследование эволюции спектральных и температурных кривых настройки второй гармоники в процессе оптического полинга германосиликатного оптического волокна излучением импульсного иттербиевого лазера.

Для достижения цели необходимо решить следующие задачи: разработать экспериментальную установку и провести оптический полинг волокна; в процессе формирования решетки на различных временных этапах измерить спектральные и температурные кривые настройки ВГГ; определить динамику их основных параметров (сдвиг пика, ширина, эффективность); выполнить численное моделирование на основе существующей теоретической модели и провести сравнительный анализ полученных результатов.



\section{Экспериментальное исследование фазового синхронизма при ГВГ в ОВ}\label{sec:nonlin2/sect1}
\clearpage

\section{Математическая модель фазового синхронизма при ГВГ в ОВ}

\section{Обсуждение результатов}

\section{Выводы к Главе 3}
\FloatBarrier